\chapter*{Introduction}
\addcontentsline{toc}{chapter}{Introduction}

%blabla

%\paragraph{test}

%\bigskip

%« Si une partie de l'historiographie anglo-saxonne récente, à l'image de Neil McLynn, a tendance à lire l'affabilité ambrosienne comme un stratagème purement politique visant à piéger l'empereur, d'autres chercheurs comme Cesare Pasini rappellent la dimension fondamentalement pastorale de l'évêque, pour qui la clémence n'est pas une arme rhétorique, mais une véritable exigence théologique. ». Bien dire que Psini voit Ambroise comme un pasteur et théologien qui agit en cohérence de sa foi, et pas un maniulateur politique moivé par ses ambitions personnelles.
